
%%%%%%%%%%  scan idx 0  |  folio cover  %%%%%%%%%%
\SGAPrintPage{cover}%
\begin{center}
{\Large Lecture Notes in Mathematics}\par
A collection of informal reports and seminars\par
Edited by A. Dold, Heidelberg and B. Eckmann, Zürich\par
\bigskip
225\par
\bigskip
Algebraic Geometry Seminar of Bois Marie 1966/67\par
SGA 6\par
Directed by\par
P. Berthelot, A. Grothendieck, L. Illusie\par
With the collaboration of\par
D. Ferrand, J.P. Jouanolou, O. Jussila,\par
S. Kleiman, M. Raynaud, J.P. Serre\par
\bigskip
{\Large Intersection Theory\par
and the Riemann--Roch Theorem}\par
\bigskip
Springer-Verlag\par
Berlin \textperiodcentered{} Heidelberg \textperiodcentered{} New York
\end{center}

%%%%%%%%%%  scan idx 1  |  folio cover-verso  %%%%%%%%%%
\SGAPrintPage{cover-verso}%
% The contemporary-template restitution is documented in the new edition preface and source records.

% The masked notice is restored from the contemporary primary template
% in Springer LNM 228 (1971), front-matter PDF page 1; see the separate
% ASTRA-FRONT-NOTICE-1971-RESTORATION evidence packet. No obscured stroke
% in the supplied SGA 6 image is claimed to have become directly readable.
This series aims to report new developments in mathematical research and teaching -- quickly, informally
and at a high level. The type of material considered for publication includes:
\begin{enumerate}
\item Preliminary drafts of original papers and monographs
\item Lectures on a new field, or presenting a new angle on a classical field
\item Seminar work-outs
\item Reports of meetings
\end{enumerate}
Texts which are out of print but still in demand may also be considered if they fall within these categories.

The timeliness of a manuscript is more important than its form, which may be unfinished or tentative.
Thus, in some instances, proofs may be merely outlined and results presented which have been or will
later be published elsewhere.

Manuscripts should comprise not less than 100 pages.

Publication of \textit{Lecture Notes} is intended as a service to the international mathematical community, in
that a commercial publisher, Springer-Verlag, can offer a wider distribution to documents which would
otherwise have a restricted readership. Once published and copyrighted, they can be documented in the
scientific literature. --

\subsection*{Manuscripts}
Manuscripts are reproduced by a photographic process; they must therefore be typed with extreme care.

Symbols not on the typewriter should be inserted by hand in indelible black ink.
Corrections to the typescript should be made by sticking the amended text over
the old one, or by obliterating errors with white correcting fluid. Authors
receive 75 free copies.

The typescript is reduced slightly in size during reproduction; best results
will not be obtained unless the text on any one page is kept within the overall
limit of $18\times26.5$ cm ($7\times10\frac12$ inches). The publishers will be
pleased to supply on request special stationery with the typing area outlined.

Manuscripts in English, German or French should be sent to Prof. Dr. A. Dold,
Mathematisches Institut der Universität Heidelberg, 69 Heidelberg/Germany,
Tiergartenstraße or Prof. Dr. B. Eckmann, Eidgenössische Technische Hochschule,
CH-8006 Zürich/Switzerland.

\subsection*{Lecture Notes in Physics}
\textbf{Bisher erschienen/Already published}

Vol. 1: J. C. Erdmann, Wärmeleitung in Kristallen, theoretische Grundlagen und
fortgeschrittene experimentelle Methoden. II, 283 Seiten. 1969. DM 20,--

Vol. 2: K. Hepp, Théorie de la renormalisation. III, 215 pages. 1969. DM 18,--

Vol. 3: A. Martin, Scattering Theory: Unitarity, Analytic and Crossing.
IV, 125 pages. 1969. DM 14,--

Vol. 4: G. Ludwig, Deutung des Begriffs physikalische Theorie und axiomatische
Grundlegung der Hilbertraumstruktur der Quantenmechanik durch Hauptsätze des
Messens. XI, 469 Seiten. 1970. DM 28,--

Vol. 5: M. Schaaf, The Reduction of the Product of Two Irreducible Unitary
Representations of the Proper Orthochronous Quantummechanical Poincaré Group.
IV, 120 pages. 1970. DM 14,--

Vol. 6: Group Representations in Mathematics and Physics. Edited by V. Bargmann.
V, 340 pages. 1970. DM 24,--

Vol. 7: R. Balescu, J. L. Lebowitz, I. Prigogine, P. Résibois, Z. W. Salsburg,
Lectures in Statistical Physics. V, 181 pages. 1971. DM 18,--

Vol. 8: Proceedings of the Second International Conference on Numerical Methods
in Fluid Dynamics. Edited by M. Holt. IX, 462 pages. 1971. DM 28,--

%%%%%%%%%%  scan idx 2  |  folio I  %%%%%%%%%%
\SGAPrintPage{I}%
\begin{center}
{\Large Lecture Notes in Mathematics}\par
A collection of informal reports and seminars\par
Edited by A. Dold, Heidelberg and B. Eckmann, Zürich\par
\bigskip
225\par
\bigskip
{\Large Intersection Theory\par
and the Riemann--Roch Theorem}\par
\bigskip
Springer-Verlag\par
Berlin \textperiodcentered{} Heidelberg \textperiodcentered{} New York 1971
\end{center}

%%%%%%%%%%  scan idx 3  |  folio II  %%%%%%%%%%
\SGAPrintPage{II}%
AMS Subject Classifications (1970): 14-02, 14C15, 14F15, 14K30, 18E30

ISBN 3-540-05647-5 Springer-Verlag Berlin \textperiodcentered{} Heidelberg \textperiodcentered{} New York\par
ISBN 0-387-05647-5 Springer-Verlag New York \textperiodcentered{} Heidelberg \textperiodcentered{} Berlin

This work is subject to copyright. All rights are reserved, whether the whole
or part of the material is concerned, specifically those of translation,
reprinting, re-use of illustrations, broadcasting, reproduction by photocopying
machine or similar means, and storage in data banks.

Under \S~54 of the German Copyright Law where copies are made for other than
private use, a fee is payable to the publisher, the amount of the fee to be
determined by agreement with the publisher.

\textcopyright{} by Springer-Verlag Berlin \textperiodcentered{} Heidelberg 1971.
Library of Congress Catalog Card Number 73-178287. Printed in Germany.

Offsetdruck: Julius Beltz, Hemsbach/Bergstr.


%%%%%%%%%%  scan idx 4  |  folio III  %%%%%%%%%%
\SGAPrintPage{III}%
\begin{center}
{\Large Algebraic Geometry Seminar of Bois Marie}\par
{\large 1966--67}\par
\vspace{1.5em}
{\Large Intersection Theory and the Riemann--Roch Theorem}\par
{\large (SGA 6)}\par
\vspace{1.5em}
a Seminar directed by\par
P. Berthelot, A. Grothendieck, L. Illusie\par
with the collaboration of\par
D. Ferrand, J.-P. Jouanolou, O. Jussila, S. Kleiman, M. Raynaud, J.-P. Serre
\end{center}

%%%%%%%%%%  scan idx 5  |  folio IV  %%%%%%%%%%
\SGAPrintPage{IV}%
%%%%%%%%%%  scan idx 6  |  folio V  %%%%%%%%%%
\SGAPrintPage{V}%
\sgaStableTarget{sga6.front.preface}\section*{\underline{Preface}}
The present edition of the Algebraic Geometry Seminar 6 reproduces, with very few changes, the first edition, distributed by the I.H.E.S. as mimeographed exposés. The corrections made concern essentially printing errors, except in Exposés I and III, where L. Illusie made slight modifications to some statements that were incorrect in the original version. In addition, footnotes have been added, in particular in Exposé XIV by A. Grothendieck, to indicate recent progress on questions that were still open at the time of the Seminar.

Finally, we point out that Exposé XI, by D. Ferrand, concerning the theory of the determinant for perfect complexes, has not been written up and will therefore not be published; to avoid errors in references, we have kept the numbering of the subsequent exposés unchanged.

\begin{flushright}June 1971\end{flushright}

%%%%%%%%%%  scan idx 7  |  folio VI  %%%%%%%%%%
\SGAPrintPage{VI}%
%%%%%%%%%%  scan idx 8  |  folio VII  %%%%%%%%%%
\SGAPrintPage{VII}%
\sgaStableTarget{sga6.front.introduction}\section*{\underline{Introduction}}
The aim of the present Seminar is to develop a \underline{global} intersection theory, and a Riemann--Roch formula, for arbitrary schemes. The reader will find a description of the Seminar's program in Exposé 0. The possibility, in principle, of proving a Riemann--Roch formula for general \underline{regular} schemes, along the lines of the well-known report of Borel--Serre in 1958, was clear from that time, at least for a projective morphism. The extension to the general case is less evident; the program of which it forms a part (and which is partially carried out in our seminar) dates from 1960. As was also the case for the duality theory of coherent sheaves (cf.\ R.~Hartshorne, Residues and Duality, Lecture Notes in Mathematics n\textsuperscript{o}~20, Springer), an essential tool for formulating a satisfactory theory is Verdier's theory of \underline{derived categories}, knowledge of which is indispensable for understanding the Seminar.

Apart from this theory, studying the Seminar requires scarcely more than a general knowledge of the foundations of Algebraic Geometry, as set out in EGA, Chapters I, II, III; in addition, we shall have only occasional references to make to EGA IV, for certain facts concerning flat or smooth morphisms, which for the most part also appear in the first exposés of SGA~1. Finally, to develop certain results with all the generality desirable for applications, we sometimes use the notions of a \underline{ringed site} and a \underline{ringed topos}, for which we refer to SGA~4, Exposés I to IV. A reader unfamiliar with the language of sites and topoi may everywhere replace these objects by ordinary topological spaces, the objects of the topos then being replaced by open subsets of these spaces; nevertheless, we advise the reader, preferably, to assimilate the language of topoi, which provides an extremely convenient principle of unification.

The fundamental notion for the theory presented here is that of a \underline{perfect complex of Modules}, and its various variants, developed in Exposés I, II, III. It seems clear that these notions, which are also important in other contexts in Algebraic Geometry (notably for formulas of Lefschetz--Verdier type in cohomology with discrete (SGA 5) or coherent coefficients), will also have their part to play outside Algebraic Geometry, notably for the formulation of a com%
%%%%%%%%%%  scan idx 9  |  folio VIII  %%%%%%%%%%
plex\SGAPrintPage{VIII}%
 analytic variant of the Riemann--Roch theorem presented here, or suitable variants of the Atiyah--Singer index theorem. Some indications in this direction will be given in Exp.\ II. Unfortunately, at the present time there is still no statement, even a heuristic one, that would encompass these two theorems (the first of which, for the moment, remains conjectural). A new idea is apparently lacking, as one is also lacking in Algebraic Geometry to obtain a proof of the Riemann--Roch theorem beyond projectivity hypotheses (cf.\ Exp.\ XIV, n\textsuperscript{o}~2).

We shall make no use in this seminar of the \underline{local theory} of intersections, set out in J.P. Serre, Algèbre Locale. Multiplicités (Lecture Notes in Mathematics n\textsuperscript{o}~11, Springer), and shall simply content ourselves with pointing out here that this course of Serre had an evident influence on the genesis of the theory in 1957. Nor shall we use the theory of the \underline{Chow ring} developed in the Chevalley Seminar 1958, Exposés 2 and 3 (École Normale Supérieure). This theory is essentially linked to nonsingularity hypotheses, whereas the aim of our seminar, on the contrary, is to develop an intersection theory on general schemes (and even on general locally ringed topoi) see on this subject the comments in Exp.\ XIV n\textsuperscript{o}~4 and n\textsuperscript{o}~8, giving the relations between our theory and that of the Chow ring, and raising the question of a generalization of the latter. We may say that the Seminar presents a coherent and self-contained intersection theory, but by no means exhausts the different points of view that are useful (or indeed indispensable) in intersection theory, and that it should consequently be supplemented by the cited exposés of Serre and Chevalley.

We have attached to the seminar (as an Appendix to Exp.\ 0) Grothendieck's 1957 report ``classes de Faisceaux et théorème de Riemann--Roch'', which served as the basis for the seminar of Borel and Serre at Princeton in the same year, as well as for their already cited article. This report outlines two proofs of the Riemann--Roch theorem for nonsingular quasi-projective varieties; the first, valid at present only in characteristic zero, but giving in return a slightly more precise result in the case of an immersion, does not yet appear in the literature (including the work of Borel--Serre and the present seminar). Reading this report of course presupposes none of the other exposés of the seminar, and may even serve as an introduction to its study on the same footing as Exposé 0, especially for a reader not yet familiar with Borel--Serre. Moreover, the proof to which we have just alluded applies essentially as it stands to the case
%%%%%%%%%%  scan idx 10  |  folio IX  %%%%%%%%%%
\SGAPrintPage{IX}%
of a projective morphism of complex analytic spaces, and will undoubtedly also apply in arbitrary characteristic once the problem of the ``exterior power'' operations in the derived category has been solved (cf.\ Exp.\ XIV, n\textsuperscript{os}~1, 2, 3). The absence of a study of this question undoubtedly constitutes the most serious gap in this Seminar, even from the point of view we have adopted.

One will notice the absence, from the table of contents of the present Seminar, of two of the participants mentioned on the cover page. One of them, J.P. Jouanolou, had taken an active part in the technical development of the first part of the seminar, but had been prevented from taking part in the oral exposés; to him we owe, in particular, the linear independence assertion contained in the important statement VI 1.1. giving the structure of the ring $K^{\bullet}$ of classes of locally free Modules on a projective bundle (which had previously been proved only under the assumption that the base scheme admits an ample invertible Module). The other, J.P. Serre, had given two oral exposés, logically independent of the rest of the Seminar, which he therefore preferred to publish as a separate article (cf.\ J.P. Serre, Groupes de Grothendieck des schémas en groupes réductifs déployés, to appear in Publications Mathématiques, n\textsuperscript{o}~34).

As in the other parts of the Algebraic Geometry Seminar of Bois Marie, the sigla SGA 1 to SGA 6 refer to the different parts of that Seminar, the Roman numeral following the siglum indicating the number of the exposé, and the Arabic numerals that follow it corresponding to the internal numbering of the exposé in question; for internal references within the present Seminar, the siglum SGA 6 is omitted, using references of the form I 4.9 (meaning: Exposé I, statement 4.9). The siglum EGA refers to the Éléments de Géométrie Algébrique of J. Dieudonné and A. Grothendieck.

%%%%%%%%%%  scan idx 11  |  folio X  %%%%%%%%%%
\SGAPrintPage{X}%
%%%%%%%%%%  scan idx 12  |  folio XI  %%%%%%%%%%
\SGAPrintPage{XI}%
\sgaStableTarget{sga6.front.table-of-contents}\section*{\underline{Table of contents}}

\noindent
\underline{EXPOSÉ 0}\\
Outline of a Program for an Intersection Theory on General Schemes\\
by A. Grothendieck \dotfill 1\\[8pt]
Classes of Sheaves and the Riemann--Roch Theorem\\
by A. Grothendieck \dotfill 20\\[8pt]
\underline{EXPOSÉ I}\\
Generalities on Finiteness Conditions in Derived Categories\\
by L. Illusie \dotfill 78\\[8pt]
\underline{EXPOSÉ II}\\
Existence of Global Resolutions\\
by L. Illusie \dotfill 160\\[8pt]
\underline{EXPOSÉ III}\\
Relative Finiteness Conditions\\
by L. Illusie \dotfill 222\\[8pt]
\underline{EXPOSÉ IV}\\
Grothendieck Groups of Ringed Topoi\\
by L. Illusie \dotfill 274\\[8pt]
\underline{EXPOSÉ V}\\
Generalities on $\lambda$-Rings\\
by P. Berthelot \dotfill 297\\[8pt]
\underline{EXPOSÉ VI}\\
The $K^{\bullet}$ of a Projective Bundle: Calculations and Consequences\\
by P. Berthelot \dotfill 365\\[8pt]
%%%%%%%%%%  scan idx 13  |  folio XII  %%%%%%%%%%
\SGAPrintPage{XII}%
\underline{EXPOSÉ VII}\\
Regular Immersions and Calculation of the $K^{\bullet}$ of a Blown-up Scheme\\
by P. Berthelot \dotfill 416\\[8pt]
\underline{EXPOSÉ VIII}\\
The Riemann--Roch Theorem\\
by P. Berthelot \dotfill 466\\[8pt]
\underline{EXPOSÉ IX}\\
Some Calculations of $K$ Groups.\\
by P. Berthelot \dotfill 498\\[8pt]
\underline{EXPOSÉ X}\\
Intersection Formalism on Proper Algebraic Schemes\\
by O. Jussila\\
With an Appendix by A. Grothendieck\\
Specialization in Intersection Theory \dotfill 519\\[8pt]
\underline{EXPOSÉ XI} -- Not written up\\[8pt]
\underline{EXPOSÉ XII}\\
A Relative Representability Theorem for the Picard Functor\\
by M. Raynaud (written up by S. Kleiman) \dotfill 595\\[8pt]
\underline{EXPOSÉ XIII}\\
Finiteness Theorems for the Picard Functor\\
by S. Kleiman \dotfill 616\\[8pt]
\underline{EXPOSÉ XIV}\\
Open Problems in Intersection Theory\\
by A. Grothendieck \dotfill 667\\[8pt]
Terminological Index \dotfill 691\\[4pt]
Index of Notations \dotfill 696\par

